\chapter{\textcolor{gray}{Analyse numérique et simulations}}
\section{Rappel de la Méthode d'Euler\cite{siciliano2012hodgkin}}
		\textbf{Problème général}\\
		Une EDO du premier ordre s'écrit sous la forme :
		\[
		\frac{dy}{dt} = f(t, y), \quad y(t_0) = y_0
		\]
		On cherche à approximer la solution $y(t)$ sur un intervalle de temps à l'aide d'un pas $\Delta t$.
		\vspace{0.5cm}
		\textbf{Formule d'Euler (mise à jour à chaque pas)}
		\[
		y_{n+1} = y_n + \Delta t \cdot f(t_n, y_n)
		\]
		où :
		\begin{itemize}
			\item $y_n$ est l'approximation de $y(t_n)$,
			\item $t_n = t_0 + n\Delta t$
		\end{itemize}
\section{Construction du problème}
	\textbf{Modèle de Hodgkin-Huxley} \\
		\begin{align*}
			\frac{dv}{dt} &= \frac{1}{C_m} \left[ I - \bar{G}_{Na} m^3 h (v - E_{Na}) - \bar{G}_K n^4 (v - E_K) - \bar{G}_l (v - E_l) \right] \\
			\frac{dn}{dt} &= \alpha_n(v)(1 - n) - \beta_n(v)n \\
			\frac{dm}{dt} &= \alpha_m(v)(1 - m) - \beta_m(v)m \\
			\frac{dh}{dt} &= \alpha_h(v)(1 - h) - \beta_h(v)h
		\end{align*}
		
		\noindent \textbf{Avec :}
		\begin{align*}
			\alpha_n(v) &= \frac{0.01(v + 50)}{1 - \exp\left( \frac{-(v + 50)}{10} \right)} 
			&\beta_n(v) &= 0.125 \exp\left( \frac{-(v + 60)}{80} \right) \\
			\alpha_m(v) &= \frac{0.1(v + 35)}{1 - \exp\left( \frac{-(v + 35)}{10} \right)} 
			&\beta_m(v) &= 4.0 \exp\left( -0.0556(v + 60) \right) \\
			\alpha_h(v) &= 0.07 \exp\left( -0.05(v + 60) \right)
			&\beta_h(v) &= \frac{1}{1 + \exp(-0.1(v + 30))}
		\end{align*}
		\vspace{0.5em}
		\noindent \textbf{Constantes :}
		\begin{align*}
			C_m &= 0.01 \, \mu\text{F/cm}^2 \\
			E_{Na} &= 55.17 \, \text{mV}, \quad
			E_K = -72.14 \, \text{mV}, \quad
			E_l = -49.42 \, \text{mV} \\
			\bar{G}_{Na} &= 1.2 \, \text{mS/cm}^2, \quad
			\bar{G}_K = 0.36 \, \text{mS/cm}^2, \quad
			\bar{G}_l = 0.003 \, \text{mS/cm}^2
		\end{align*}
	
\section{Implémentation}
\subsection{L'algorithme}
\begin{figure}[h]
	\centering
	\begin{minipage}{0.45\textwidth}
		\begin{itemize}
			\item \textbf{Démarrer}
			\item[] \vspace{0.5em}
			\item \textbf{Initialisation}
			\item[] \vspace{0.5em}
			\item \textbf{Boucle : $i = 1$ à $N-1$}
			\item[] \vspace{0.5em}
			\item \textbf{Appel fonction HH}
			\item[] \vspace{0.5em}
			\item \textbf{Mise à jour}
			\item[] \vspace{0.5em}
			\item \textbf{Décision : $i < N-1$}
			\item[] \vspace{0.5em}
			\item \textbf{Tracer les courbes}
			\item[] \vspace{0.5em}
			\item \textbf{Fin}
		\end{itemize}
	\end{minipage}%
	\hfill
	\begin{minipage}{0.5\textwidth}
		\centering
		\includegraphics[scale=0.3]{pictures/schema.png}
	\end{minipage}
	\caption{Schéma de l'algorithme avec le modèle Hodgkin-Huxley}
	\label{fig:schema_hh}
\end{figure}
\subsection{Les fonctions}
		\begin{center}
			\begin{minipage}{0.49\textwidth}
				\begin{lstlisting}
				% Fonctions des variables de porte 
				function a = am(V)
				a = 0.1 * (V + 35) / (1 - exp(-(V + 35) / 10));
				end
				function b = bm(V)
				b = 4.0 * exp(-0.0556 * (V + 60));
				end
				function a = an(V)
				a = 0.01 * (V + 50) / (1 - exp(-(V + 50) / 10));
				end
				\end{lstlisting}
			\end{minipage}%
			\hfill
			\begin{minipage}{0.49\textwidth}
				\begin{lstlisting}
				function b = bn(V)
				b = 0.125 * exp(-(V + 60) / 80);
				end
				function a = ah(V)
				a = 0.07 * exp(-0.05 * (V + 60));
				end
				function b = bh(V)
				b = 1 / (1 + exp(-0.1 * (V + 30)));
				end
				\end{lstlisting}
			\end{minipage}
		\end{center}
\subsection{La Fonction HH}
	\begin{center}
		\begin{minipage}{0.49\textwidth}
			\begin{lstlisting}
			% Fonctions utilisees
			function dydt = HH(t, y)
			% Constantes
			ENa = 55.17;    % mV
			EK = -72.14;    % mV
			El = -49.42;    % mV
			gbarNa = 1.2;   % mS/cm^2
			gbarK = 0.36;   % mS/cm^2
			gbarl = 0.003;  % mS/cm^2
			Cm = 0.01;      % uF/cm^2
			I = 0.1;        % uA/cm^2
			% Variables d'etat
			V = y(1);
			n = y(2);
			m = y(3);
			h = y(4);
			% Conductances ioniques
			gNa = gbarNa * m^3 * h;
			gK = gbarK * n^4;
			gl = gbarl;
			\end{lstlisting}
		\end{minipage}%
		\hfill
		\begin{minipage}{0.49\textwidth}
			\begin{lstlisting}
			% Courants ioniques
			INa = gNa * (V - ENa);
			IK = gK * (V - EK);
			Il = gl * (V - El);
			% Equations differentielles
			dVdt = (1 / Cm) * (I - INa - IK - Il);
			dndt = an(V)*(1 - n) - bn(V)*n;
			dmdt = am(V)*(1 - m) - bm(V)*m;
			dhdt = ah(V)*(1 - h) - bh(V)*h;
			dydt = [dVdt; dndt; dmdt; dhdt];
			end
			\end{lstlisting}
		\end{minipage}
	\end{center}
\subsection{Script}
		\begin{center}
			\begin{minipage}{0.49\textwidth}
				\begin{lstlisting}
				% Script principal 
				T = 30;
				pas = 0.0001;
				H = 0:pas:T;
				y0 = [-65, 0.3177, 0.0529, 0.5961]'; % Conditions initiales
				Y = [y0 zeros(4,length(H)-1)];
				for i = 1:length(H)-1
				Y(:,i+1) = Y(:,i) + pas * HH(H(i), Y(:,i));
				end
				subplot(1,2,1)
				plot(H, Y(1,:));
				xlabel('Temps (ms)');
				ylabel('V (mV)');
				title('Potentiel membranaire');
				grid on;
				\end{lstlisting}
			\end{minipage}%
			\hfill
			\begin{minipage}{0.49\textwidth}
				\begin{lstlisting}
				subplot(1,2,2)
				hold on
				a2=Y(2,:);
				a3=Y(3,:);
				a4=Y(4,:);
				z = T/(pas*6);
				plot(H(1:z), a2(1:z), 'r', 'DisplayName', 'n');
				plot(H(1:z), a3(1:z), 'g', 'DisplayName', 'm');
				plot(H(1:z), a4(1:z), 'b', 'DisplayName', 'h');
				xlabel('Temps (ms)');
				ylabel('Variables de porte');
				title('n, m, h dans le modele de HH');
				legend show
				grid on;
				hold off
				\end{lstlisting}
			\end{minipage}
		\end{center}
\section{Résultats}
	\subsection{La solution obtenue avec la méthode d'Euler}
	\begin{figure}[h]
		\centering
		\includegraphics[width=\linewidth]{pictures/simulation.png}
		\caption{La solution obtenue avec la méthode d'Euler}
	\end{figure}
\textbf{\textit{Interprétation du graphe I :}}
Le graphique de gauche représente l'évolution du potentiel membranaire d'un neurone selon le modèle de Hodgkin-Huxley. On y observe deux potentiels d'action successifs, caractérisés par des pics d'activité électrique. Le potentiel débute à environ -65 mV, correspondant à l'état de repos de la membrane, puis monte rapidement vers +45 mV lors de la phase de dépolarisation. Ensuite, il chute en dessous du potentiel de repos (hyperpolarisation) avant de revenir à l'état initial. Chaque pic traduit une réponse du neurone à une stimulation suffisante, déclenchant un influx nerveux. Ce comportement résulte de l'activation coordonnée des canaux ioniques, et reflète précisément la dynamique électrique neuronale modélisée par les équations de Hodgkin-Huxley.\\
\textbf{\textit{Interprétation du graphe II :}}
Le graphique de droite illustre l'évolution des variables de porte \( m \), \( h \) et \( n \) dans le modèle de Hodgkin-Huxley, qui décrivent respectivement l'activation et l'inactivation des canaux sodiques, ainsi que l'activation des canaux potassiques. La variable \( m \) (courbe verte) augmente très rapidement dès le début de la stimulation, traduisant l'ouverture rapide des canaux sodium, ce qui permet l'entrée massive de $Na^{+}$ dans la cellule et provoque la dépolarisation de la membrane. Parallèlement, la variable \( h \) (courbe bleue) diminue progressivement, représentant l'inactivation des canaux sodium, processus qui contribue à l'arrêt du flux de $Na^{+}$ et à la fin du potentiel d'action. Quant à la variable \( n \) (courbe rouge), elle augmente plus lentement, reflétant l'ouverture progressive des canaux potassium, permettant la sortie de $K^{+}$ et participant ainsi à la repolarisation de la membrane, puis à l'hyperpolarisation. L'évolution de ces trois variables rend compte avec précision de la régulation ionique à l'origine du potentiel d'action neuronal.

\subsection{La solution avec différentes Méthodes}
\begin{figure}[h]
	\centering
		\includegraphics[width=\linewidth]{pictures/sumi-diff.png}
	\caption{Comparaison entre différentes méthodes numériques}
\end{figure}

